\documentclass[12pt,a4paper]{article}

% Pacotes essenciais
\usepackage[utf8]{inputenc}
\usepackage[T1]{fontenc}
\usepackage[brazil]{babel}
\usepackage[top=3cm, bottom=2cm, left=3cm, right=2cm]{geometry}
\usepackage{setspace}
\usepackage{indentfirst}
\usepackage{graphicx}
\usepackage{titlesec}
\usepackage{booktabs}
\usepackage{amsmath}

% Configuração de espaçamento (1.5)
\setstretch{1.5}

\begin{document}

% ---------------------------------------------------
% CAPA
% ---------------------------------------------------
\begin{titlepage}
    \begin{center}
        {\large \textbf{Victor Hugo Bitencourt}} \\[7cm]
        
        {\LARGE \textbf{Arquitetura de Baixo Custo baseada em Redes Sem Fio para Monitoramento Centralizado de Experimentos Educacionais em Laboratórios de Física}} \\[8cm]
        
        {\large Salvador-BA} \\
        {\large 28 de maio de 2026}
    \end{center}
\end{titlepage}

% ---------------------------------------------------
% FOLHA DE ROSTO
% ---------------------------------------------------
\begin{titlepage}
    \begin{center}
        {\large \textbf{Victor Hugo Bitencourt}} \\[5cm]
        
        {\LARGE \textbf{Arquitetura de Baixo Custo baseada em Redes Sem Fio para Monitoramento Centralizado de Experimentos Educacionais em Laboratórios de Física}} \\[3cm]
        
        \begin{flushright}
            \begin{minipage}{8cm}
                Anteprojeto de Trabalho de Conclusão de Curso apresentado ao curso de Bacharelado em Sistemas de Informação.\\
                \textbf{Área da Computação:} Sistemas Embarcados / Internet das Coisas (IoT)\\
                \textbf{Orientador:} Prof. Dr. Cláudio Alves de Amorim
            \end{minipage}
        \end{flushright}
        
        \vspace{4cm}
        
        \textsc{UNEB - Universidade do Estado da Bahia}\\
        \textsc{Departamento de Ciências Exatas e da Terra} \\[2cm]
        
        {\large Salvador-BA} \\
        {\large 28 de maio de 2026}
    \end{center}
\end{titlepage}

% ---------------------------------------------------
% FOLHA DE APROVAÇÃO
% ---------------------------------------------------
\begin{titlepage}
    \begin{center}
        {\Large \textbf{Folha de Aprovação}} \\[2cm]
    \end{center}
    
    Anteprojeto sob o título \textit{Arquitetura de Baixo Custo baseada em Redes Sem Fio para Monitoramento Centralizado de Experimentos Educacionais em Laboratórios de Física} apresentado como exigência parcial para avaliação na disciplina Trabalho de Conclusão de Curso I do bacharelado em Sistemas de Informação da Universidade do Estado da Bahia entregue por \textit{Victor Hugo Bitencourt} a Dr. Jorge Alberto Prado de Campos, professor da disciplina, em 2 de junho de 2026, em Salvador, Bahia.
    
    \vspace{3cm}
    
    \begin{center}
        \rule{10cm}{0.1mm} \\
        Victor Hugo Bitencourt \\
        Orientando \\[2cm]
        
        \rule{10cm}{0.1mm} \\
        Prof. Dr. Cláudio Alves de Amorim \\
        Orientador
    \end{center}
\end{titlepage}

% ---------------------------------------------------
% SUMÁRIO
% ---------------------------------------------------
\newpage
\tableofcontents
\thispagestyle{empty}
\newpage

\setcounter{page}{4} % Ajusta a numeração da página para bater com o sumário

% ---------------------------------------------------
% 1. INTRODUÇÃO
% ---------------------------------------------------
\section{Introdução}

O ensino experimental em laboratórios de Física é fundamental para a consolidação teórica dos estudantes, permitindo a observação empírica de fenômenos modelados em sala de aula. No entanto, a instrumentação comercial utilizada para a coleta e análise de dados baseia-se frequentemente em sistemas proprietários de aquisição de dados (DAQ). Esses equipamentos possuem alto custo de aquisição e manutenção, o que afeta diretamente a realidade econômica brasileira, engessando a escalabilidade e dificultando a modernização da infraestrutura em instituições públicas de ensino.

Conforme identificado na revisão sistemática da literatura previamente conduzida neste escopo, há uma tendência crescente na adoção de abordagens baseadas em sistemas embarcados que buscam alternativas viáveis de baixo custo. Embora a literatura contemple inovações, nota-se que o uso de ferramentas alternativas e não padronizadas para baratear os experimentos --- a exemplo da utilização de smartphones pessoais de estudantes para coleta de dados --- pode dificultar a gestão pedagógica, gerando distrações e falta de padronização nas aulas. Além disso, as soluções que empregam microcontroladores frequentemente operam de forma isolada, exigindo conexão física (via interface USB) a computadores dedicados por bancada. Essa exigência arquitetural gera ociosidade de hardware e inviabiliza o monitoramento centralizado dos resultados, dificultando o acompanhamento simultâneo de múltiplos grupos pelo docente.

Diante desse cenário, surge a necessidade de explorar soluções que combinem a acessibilidade financeira dos microcontroladores à flexibilidade das redes sem fio, com foco primordial na aplicabilidade educacional e na facilidade de implementação, garantindo uma coleta de dados fluida, padronizada e integrada. 

Desta forma, o presente trabalho define como problema de pesquisa: \textbf{Como desenvolver e validar uma arquitetura de baixo custo, baseada em comunicação sem fio, focada na usabilidade educacional para a aquisição e monitoramento centralizado de experimentos em laboratórios de Física?}

% ---------------------------------------------------
% 2. OBJETIVOS
% ---------------------------------------------------
\section{Objetivos}

\subsection{Objetivo Geral}
Construir e validar um sistema de aquisição, monitoramento e controle para múltiplos experimentos didáticos de Física, utilizando microcontroladores de baixo custo integrados a um servidor centralizado via rede sem fio (Wi-Fi), com acesso por interface gráfica.

\subsection{Objetivos Específicos}
\begin{itemize}
    \item Avaliar plataformas microcontroladas para estabelecer a relação ideal entre custo, processamento, capacidade de leitura e conectividade sem fio.
    \item Projetar a arquitetura de software (aplicação web/servidor e nós embarcados), selecionando metodologias de processamento adequadas para o recebimento de telemetria centralizada.
    \item Implementar a arquitetura desenvolvida em uma Prova de Conceito (PoC) utilizando um arranjo experimental didático de fácil montagem e alta relevância pedagógica.
    \item Validar o sistema aferindo seu custo de implantação e a qualidade da experiência de uso, com foco na adequação da interface e viabilidade no contexto educacional.
\end{itemize}

% ---------------------------------------------------
% 3. JUSTIFICATIVAS E CONTRIBUIÇÕES
% ---------------------------------------------------
\section{Justificativas e Contribuições}

A motivação central deste trabalho é educacional e econômica. Atualmente, os dispositivos experimentais didáticos funcionam sem conexão ou, quando conectados, exigem hardware dedicado que permanece ocioso grande parte do tempo. O desenvolvimento de uma arquitetura centralizada visa minimizar essa ociosidade, permitindo que laboratórios escolares e universitários sejam equipados a um custo substancialmente reduzido, operando múltiplos experimentos simultâneos orquestrados por um único ponto central de processamento.

No âmbito acadêmico, a revisão sistemática apontou para uma lacuna pertinente: embora existam muitos projetos focados na precisão, poucos buscam equilibrar de fato o custo da arquitetura e a robustez técnica com a adequação pedagógica para professores e alunos, limitando a adoção real dessas tecnologias nas instituições. O uso de uma plataforma padronizada substitui improvisos e propicia uma melhor condução do experimento.

A contribuição tecnológica assegura-se pela abordagem pragmática da infraestrutura integrada: fornecer um ambiente em que o orientador do laboratório possa gerenciar múltiplos nós sensores, persistir e analisar dados experimentais através de uma interface gráfica amigável, livre de emaranhados de cabos e dependências comerciais limitantes. Isso abre portas também para trabalhos futuros no âmbito de experimentos conduzidos de forma remota ou híbrida.

% ---------------------------------------------------
% 4. METODOLOGIA
% ---------------------------------------------------
\section{Metodologia}

A presente pesquisa classifica-se, quanto à sua natureza, como \textbf{Aplicada}, pois busca resolver um problema prático do ambiente de ensino. Possui objetivos de caráter \textbf{Exploratório}, devido à busca por familiaridade, identificação de lacunas e avaliação de alternativas de hardware. Possui abordagem \textbf{Quantitativa e Qualitativa}, por tratar de métricas de software e medição de custos associadas à análise qualitativa da experiência educacional Configura-se proceduralmente como um \textbf{Desenvolvimento Tecnológico}, por culminar na entrega de um sistema funcional completo. O trabalho será conduzido nas seguintes etapas:

\textbf{1. Especificação de Hardware e Software:} Realizar-se-á uma avaliação das tecnologias embarcadas (como o ecossistema ESP32 com Wi-Fi integrado ou alternativas como o Arduino Uno R3 em conjunto com módulos externos de conectividade) considerando suas limitações e vantagens de custo. No software, a \textit{stack} tecnológica será delimitada com foco na comunicação eficiente entre os microcontroladores e a interface gráfica baseada em web (alojada, por exemplo, em um Raspberry Pi atuando como servidor/banco de dados central).

\textbf{2. Desenvolvimento da Arquitetura:} Implementação do \textit{firmware} responsável por realizar a coleta nos nós experimentais e enviar a telemetria; em paralelo, ocorrerá a implementação do sistema centralizado para persistência e visualização dos resultados na interface gráfica.

\textbf{3. Montagem da Prova de Conceito (Experimento):} Visando facilitar a adoção e a montagem pelo corpo docente, a Prova de Conceito dará prioridade a arranjos de fácil execução e com forte base pedagógica (suportados pelos resultados da revisão sistemática). Como prioridade, utilizar-se-á o experimento de Eletromagnetismo (Circuitos) de \textbf{Carga e Descarga de Capacitor}. Secundariamente, havendo necessidade, poderão ser testados ensaios de Cinemática, como a \textbf{Queda Livre}. A definição exata será consolidada durante o ciclo de projeto.

\textbf{4. Validação e Coleta de Dados:} O foco primário da validação observará os aspectos práticos:
\begin{itemize}
    \item \textit{Custo da Arquitetura:} A análise será mediada por uma comparação sistemática do custo total de aquisição entre a solução proposta e as abordagens comerciais existentes, bem como aquelas levantadas durante a revisão da literatura.
    \item \textit{Avaliação de Uaabilidade} A análise será mediada por uma avaliação qualitativa via inspeção de usabilidade e/ou eficácia pedagógica, sem usuários finais.
\end{itemize}

% ---------------------------------------------------
% 5. CRONOGRAMA
% ---------------------------------------------------
\section{Cronograma}

As atividades relacionadas ao desenvolvimento e execução da segunda etapa do Trabalho de Conclusão de Curso (TCC II) estão planejadas para o segundo semestre letivo de 2026, conforme disposto na Tabela 1.

\begin{table}[h]
\centering
\caption{Cronograma de Execução - TCC II (Julho a Dezembro de 2026)}
\vspace{0.3cm}
\begin{tabular}{lcccccc}
\toprule
\textbf{Atividade} & \textbf{Jul} & \textbf{Ago} & \textbf{Set} & \textbf{Out} & \textbf{Nov} & \textbf{Dez} \\
\midrule
Definição da plataforma de HW e experimento & $\bullet$ & & & & & \\
Desenvolvimento do Firmware (Nós sensores) & $\bullet$ & $\bullet$ & & & & \\
Desenvolvimento da Aplicação Web e BD      & & $\bullet$ & $\bullet$ & & & \\
Montagem e integração da Prova de Conceito & & & $\bullet$ & $\bullet$ & & \\
Testes finais de comportamento do sistema e validação & & & & $\bullet$ & & \\
Análise e consolidação dos resultados      & & & & $\bullet$ & $\bullet$ & \\
Redação final da monografia                & & & & & $\bullet$ & $\bullet$ \\
Defesa do Trabalho de Conclusão de Curso   & & & & & & $\bullet$ \\
\bottomrule
\end{tabular}
\end{table}

% ---------------------------------------------------
% REFERÊNCIAS
% ---------------------------------------------------
\newpage
\renewcommand\refname{Referências}
\begin{thebibliography}{9}

\bibitem{mica2012}
HEMOND, B. D. et al. MICA: an innovative approach to remote data acquisition. \textit{IEEE Instrumentation \& Measurement Magazine}, v. 15, n. 5, p. 22-27, 2012.

\bibitem{essaadaoui2021}
ESSAADAOUI, R. et al. Construction of an educational device for real time data acquisition based on arduino for a calorimetric study. In: \textit{2021 9th International Conference on Information and Education Technology (ICIET)}. IEEE, 2021. p. 21-25.

\bibitem{santos2026}
SANTOS, J. F. M. d.; NUNES, L. A. d. O. Integrating microcontrollers and data analysis in mechanics education: A stem approach to free fall. \textit{IEEE Revista Iberoamericana de Tecnologias del Aprendizaje}, v. 21, p. 101-107, 2026.

\end{thebibliography}

\end{document}
